\section{Límites de la evidencia sobre la emergence y sus implicaciones de ingeniería}

En este punto, las diapositivas dejan de añadir casos y pasan a resumir qué puede respaldar la emergence y qué no. Puede mostrar que seguir escalando los modelos quizá abra regiones de tareas que los experimentos con modelos pequeños no pueden predecir. También puede recordar a la investigación en seguridad que preste atención a posibles undesirable capabilities. Sin embargo, todavía no permite predecir umbrales precisos, describir el límite a escala infinita ni establecer la proporción óptima de la arquitectura.

\lecturefigure{slide-19.jpg}{Resumen de la emergence: escala, datos, fine-tuning y preguntas abiertas}{Diapositiva oficial 19; Wei et al., 2022}

\begin{knowledgebox}{Lectura de la figura: separar el resumen en observaciones, inferencias y preguntas abiertas}
La observación es que varias few-shot abilities y prompting techniques muestran un comportamiento de umbral en las familias de modelos evaluadas. La inferencia de ingeniería es que el límite de un método no debe juzgarse solo con modelos pequeños. Al mismo tiempo, conviene explorar cómo los datos y el posentrenamiento pueden desplazar ese límite hacia la izquierda. Las preguntas abiertas incluyen cómo predecir capacidades nuevas, cómo elegir el eje horizontal correcto, cómo hacer ablaciones confiables con el costo de un preentrenamiento a gran escala y si las undesirable abilities también podrían aparecer de forma similar.
\end{knowledgebox}

\teachervoice{El ponente llama a la emergence un implicit argument para seguir con el scaling, porque una expansión costosa puede traer capacidades difíciles de enumerar de antemano. También afirma con claridad que predecir emergencias futuras sigue siendo difícil y que no conoce un método suficientemente general. No se trata de sostener que «la escala lo resuelve todo», sino de un argumento de investigación sobre beneficios y riesgos desconocidos.}

\subsection{Por qué el preentrenamiento a gran escala difícilmente produce una teoría limpia}

Un estudiante pregunta cómo deberían repartirse los nuevos parámetros entre depth, attention heads y embedding width. El ponente responde que en ese momento no existía una principled recipe; muchas decisiones dependían de la capacidad de las TPU y de la experiencia de ingeniería. Además, una pretraining ablation completa es extremadamente costosa. Por eso suelen verse las curve de unas pocas model family, pero falta el control de una sola variable que exigiría un diseño experimental ideal.

\begin{warningbox}{«No haber hecho ablaciones» no es «un problema menor»}
Si cambian al mismo tiempo data, tokenizer, architecture y optimizer, el desplazamiento de un umbral solo puede atribuirse a la recipe completa. Reducirlo a «más parámetros hacen que aparezca la capacidad» es una atribución excesiva. Los apuntes, las revisiones de artículos y los comunicados de producto deben distinguir entre las relaciones observables y los mecanismos no verificados.
\end{warningbox}

\subsection{Los Benchmark también envejecen}

Algunas tareas de BIG-Bench se mantuvieron flat con GPT-3/LaMDA. Cuando apareció PaLM, decenas de ellas superaron el baseline. Incluso durante la clase se descubrió que ChatGPT quizá ya había resuelto una tarea de la lista. Por eso, un Benchmark no es una escala estática de dificultad. Las mejoras de los modelos, los métodos de prompting y la contaminación de datos pueden hacer que una etiqueta «hard» pierda validez rápidamente.

\begin{importantbox}{Una lista reutilizable para auditar afirmaciones sobre emergence}
\begin{enumerate}
  \item ¿Se reporta un baseline de chance, human o prior-system?
  \item ¿Se cubren model scales suficientemente densas y se reportan barras de error o varias seed?
  \item ¿Cambian al mismo tiempo data, architecture y post-training entre familias de modelos?
  \item ¿La metric es continua o convierte una calidad continua en un umbral mediante exact match?
  \item ¿Se pueden reproducir los resultados con un prompt nuevo, otra task split y una familia de modelos independiente?
\end{enumerate}
\end{importantbox}

\subsection{Resumen de la sección}

La emergence es un empirical pattern que vale la pena medir, pero no es una teoría causal completa. Su uso más valioso es mostrar dónde falla la extrapolación. Un método ineficaz en modelos pequeños puede funcionar en modelos grandes, una tendencia negativa que parece estable en modelos pequeños puede invertirse y los benchmark antiguos también pueden superarse rápidamente.
